\uuid{DKk0}
\titre{Étude de fonction}
\auteur{bodin ; hanani ; mzari}
\organisation{exo7}
\licence{CC BY-NC-SA 4.0}
\datecreate{2020-01-20}
\niveau{L1}
\module{Analyse}
\chapitre{Fonctions circulaires et hyperboliques inverses}
\sousChapitre{Autre}
\difficulte{3}
\theme{Fonctions usuelles, étude de fonction}
\exo7qcmid{lille-1.fonctions-usuelles.028}

\contenu{
\begin{question}
\qclassification{126.00}
\qtags{motcle=étude de fonction}
Soit $f(x)= \exp ( \frac{\ln^2 |x|}{\ln^2 |x|+1})$. On notera $D_f$ le domaine de définition de $f$.  Quelles sont les assertions vraies ?
\begin{answers}
    \bad{$D_f=]0,+\infty[$}

     \good{$f$ est paire.}

     \bad{$f$ est croissante sur $]0,+\infty[$.}

    \good{$f$ est une bijection de $]0,1]$ dans $[1,e[$.}

\end{answers}
\begin{explanations}
$D_f= \Rr^*$. Pour tout $x\in \Rr^*$, $f(-x)=f(x)$, donc $f$ est paire.

$f$ est dérivable sur $\Rr^*$. Si $x>0$,  $f'(x)= \frac{2\ln x}{x(\ln^2x+1)^2}f(x)$. En étudiant les variations de $f$, on déduit que $f$ est n'est pas monotone sur $]0,+\infty[$ et qu'elle établit une bijection de $]0,1]$ dans $[1,e[$ et de $[1,+\infty[$ dans $[1,e[$.

\end{explanations}
\end{question}
}
